\subsection{Режим «Обучение»: книга и симулятор вместе}
\label{sec:learning-mode}

В симуляторе есть режим \ui{Обучение}. Он ведёт по тем же темам, что и книга, но в виде коротких \textbf{миссий}: читаете раздел книги, делаете задание в рамке, затем закрепляете его миссией в симуляторе. Это один маршрут, а не два разных курса.

\begin{shotcallouts}{10-mission-step.png}
    \callout{0.965}{0.915}{1}
    \callout{0.80}{0.885}{2}
    \callout{0.885}{0.885}{3}
    \callout{0.85}{0.795}{4}
    \callout{0.795}{0.725}{5}
    \callout{0.86}{0.665}{6}
    \callout{0.82}{0.06}{7}
\end{shotcallouts}

\begin{enumerate}[label=\bubble{\arabic*}]
    \item Сколько миссий пройдено из общего числа выбранного языка (на рисунке 0 из 72).
    \item Язык миссий: \ui{Blockly} (16 миссий), \ui{Lua} (72) или \ui{Python} (108). У каждого языка своя карта и свой прогресс.
    \item Что показывать на экране: \ui{Код} (редактор) или \ui{Сцену} (чтобы наблюдать за дроном без редактора).
    \item Текущий шаг миссии. Шаги короткие: один шаг --- одно действие.
    \item \ui{Сделал, дальше} переходит к следующему шагу; \ui{Подсказка} раскрывает намёк, а в раскрытой подсказке есть кнопка \ui{Показать решение}: используйте её, когда застряли, и обязательно разберите, чем решение отличается от вашего.
    \item \ui{Карта миссий}: все миссии выбранного языка по главам.
    \item \ui{Выйти из обучения} возвращает свободный симулятор. Написанная программа остаётся в редакторе.
\end{enumerate}

\subsubsection{Как устроена миссия}

Миссия проходит по порядку: вводная карточка с целью, затем шаги, затем кнопка \ui{Проверить полёт}. Проверка --- не «галочка за старание». Симулятор смотрит на журнал настоящего выполнения: какие команды приняты, какие состояния прошёл автопилот, где дрон был и что напечатала программа. Маршрут засчитывается по порядку фактических прибытий в точки, посадка --- по состоянию автопилота и высоте. Поэтому программа, которая только печатает «готово», миссию не пройдёт.

Если проверка не прошла, симулятор объясняет причину человеческим языком и показывает, что именно не совпало. Исправьте программу и нажмите \ui{Проверить полёт} снова.

\subsubsection{Карта миссий и порядок прохождения}

\begin{center}
\simshot[0.9\linewidth]{09-mission-map.png}
\captionof{figure}{Карта миссий для Lua. Следующая миссия открывается только после предыдущей; закрытые отмечены замком.}
\end{center}

Миссии сгруппированы в главы. На карте это блоки по девять миссий: \textit{Дрон и первая команда}, \textit{Сигналы и подсветка}, \textit{Маленькие программы}, \textit{Время и события} и так далее. Понятие вводится до того, как им пользуются: сначала подготовка двигателей на земле, потом свет, потом переменные и условия, потом время, и только затем полёт.

\subsubsection{Аккаунт и сохранение прогресса}

Панель \ui{Аккаунт} необязательна. Без входа прогресс хранится в браузере этого компьютера. Если нужно, чтобы прогресс не терялся и был виден преподавателю, войдите в аккаунт: прогресс по трём языкам хранится на сервере. Сессия действует два часа без обращений к серверу, после этого нужно войти заново. Учитель видит список своих учеников и число пройденных миссий по каждому языку.

\subsubsection{Как книга соотносится с миссиями}

Вверху каждой главы книги стоит зелёная врезка \faGamepad\ \textit{Миссии симулятора} с номерами миссий по теме главы, а после урока --- рамка \faPencilAlt\ \textit{Задание}. Сначала решаем задание в рамке своими руками, затем проходим миссию: в ней те же идеи, но с автоматической проверкой.

\begin{center}
\small
\begin{tabularx}{\linewidth}{@{}>{\bfseries}X X X@{}}
\toprule
Глава книги & Миссии Lua & Миссии Python \\
\midrule
Подготовка (гл.~\ref{sec:start}) & 1--9 & 1--9 \\
Основы языка (Lua: гл.~\ref{sec:lua}; Python: часть~II) & 19--27 & 19--27 \\
Светодиоды и таймеры (гл.~\ref{sec:leds}) & 10--18, 28--36 & 10--18, 28--36 \\
Автономные миссии (гл.~\ref{sec:missions}) & 37--72 & 37--72 \\
Камера, ArUco и наведение & --- & 73--99 \\
Самостоятельные миссии & --- & 100--108 \\
\bottomrule
\end{tabularx}
\end{center}

Камера и ArUco изучаются только на Python: у Lua-автопилота Пионера нет API камеры. Блочный практикум (16 миссий) лучше пройти параллельно с первыми главами: он показывает те же команды в виде блоков.

\begin{rulebox}[Порядок работы]
Читаем раздел $\to$ запускаем примеры в редакторе $\to$ решаем задание в рамке $\to$ проходим миссию в панели \ui{Обучение}. Если миссия не проходится, вернитесь к разделу: в нём объяснено, зачем нужен каждый шаг.
\end{rulebox}
